\chapter{Preface}

This book is a rigorous scientific and architectural monograph on procedural audio processing, synthesizer construction, and deterministic Digital Audio Workstation (DAW) engineering. Its concrete reference vehicle is \textbf{B-Sound}, the Virtual Digital Signal Processor (Virtual DSP) of the B-System / BTRON~3.20 real-time operating system, together with its canonical stage DAW \textbf{Orchestra} and its phase-locked audio-visual runtime \textbf{Lil~Cu~64}.

Modern audio software is routinely split across three disconnected worlds. Music theorists study algebraic pitch structures, group actions on triads, and metric hierarchies in isolation from hardware. Digital signal processing (DSP) textbooks derive discrete-time filters, frequency modulation, and reverberators as abstract difference equations without addressing real-time memory budgets or scheduling deadlines. Systems programmers build audio engines and DAWs around gigabyte sample libraries, dynamic plugin loaders, and best-effort desktop threads whose worst-case execution time cannot be bounded.

\textbf{B-Sound} collapses those three worlds into one verifiable mathematical and engineering artifact. Every musical pitch belongs to an exact equal-tempered cyclic group $\mathbb{Z}_{12}$ mapped onto a bounded lookup table. Every instrument tone is a stateless, zero-byte procedural DSP procedure—a \emph{Sound Shim}—built from closed-form oscillators, biquad filters, Chebyshev wave-shapers, and Schroeder comb-allpass reverberators. Every performance runs inside a statically allocated, nine-element C99 signal graph whose worst-case execution time ($632\,\mu\text{s}$ per $Q=32$ sample quantum) is provably below the hardware audio period ($666\,\mu\text{s}$ at $48\,\text{kHz}$, $725\,\mu\text{s}$ at $44.1\,\text{kHz}$).

The text is written for graduate students, DSP engineers, synthesizer and emulator builders, and real-time systems architects. Part~I establishes the mathematical substrate: the Lie-group and dihedral-group structure of musical harmony, followed by analytical waveform synthesis, FM Bessel spectra, biquad filter design, non-linear saturation, and artificial reverberation. Part~II formalizes the B-Sound Virtual DSP, its 23 canonical Sound Shims, its 32-program registry, its loudness-decay voice stealer, and its phase-locked coupling to 4D Hopf fibration geometry and critically damped 3D kinematics. Part~III dissects the reference RTOS DAW \emph{Orchestra} from first principles: three-layer isolation, integer musical time, static graph admission control, shared-hall memory pools, and deterministic binary containers (73,408-byte sessions, Standard MIDI Files~1.0, and write-only RIFF/WAVE export). Part~IV demonstrates the complete system through its four-track procedural album and situates B-Sound among classical synthesizers, hardware emulators, and commercial DAWs.

Every constant printed in this book—$48000\,\text{Hz}$, $Q=32$, $632\,\mu\text{s}$, $48$ voice slots, $32$ programs, $960$ ticks per quarter note, and $73408$ session bytes—is a literal constant enforced by the running B-System implementation. Where the architecture omits a feature or carries a documented limit, the book states the exact invariant that closes it.
